\section*{الفصل \ref{ch:b3:sequence-stratigraphy} --- الستراتيغرافيا التتابعية}

\begin{solution}{exo:b3:sequence-stratigraphy:1}
$A = 0.5 + 0.1 = \qty{0.6}{m/kyr}$؛ ففي \qty{20}{kyr}،
\qty{12}{m} من حيّز الترسيب. ومع انخفاض مستوى البحر، $A = -0.3 + 0.1 = \qty{-0.2}{m/kyr}$: انخفاض
نسبيّ يتباطأ بفعل الهبوط.
\end{solution}

\begin{solution}{exo:b3:sequence-stratigraphy:2}
(أ) \omterm{def:b3:sequence-stratigraphy:tract}{مسلك أنظمة طور الانخفاض} (انحسار قسري). (ب) \omterm{def:b3:sequence-stratigraphy:tract}{مسلك الأنظمة التجاوزي}. (ج)
\omterm{def:b3:sequence-stratigraphy:tract}{مسلك أنظمة المستوى المرتفع}. (د) \omterm{def:b3:sequence-stratigraphy:tract}{مسلك أنظمة المستوى المنخفض} (انحسار عادي بعد الانخفاض).
\end{solution}

\begin{solution}{exo:b3:sequence-stratigraphy:3}
تنتهي كلّ متتالية فرعية أقلّ عمقًا من التي تحتها: تراكم تقدّمي،
يتحرّك خط الشاطئ نحو البحر، كما في \omterm{def:b3:sequence-stratigraphy:tract}{مسلك أنظمة المستوى المرتفع} أو المنخفض. أمّا
الترتيب المعاكس، حيث تنتهي كلّ متتالية أدقّ وأعمق، فهو تراجعيّ:
\omterm{def:b3:sequence-stratigraphy:tract}{مسلك الأنظمة التجاوزي}.
\end{solution}

\begin{solution}{exo:b3:sequence-stratigraphy:4}
خمسة أزواج لكلّ حزمة هي نسبة الانحراف المركزي القصير (نحو
\qty{100}{kyr}) إلى التقدّم الاعتدالي (نحو \qty{20}{kyr}): كلّ زوج هو دورة تقدّم،
وكلّ حزمة دورة انحراف مركزي. ويكون المعدّل
$1.6/100 = \qty{1.6}{cm/kyr}$، ويدوم زوج بسمك \qty{0.32}{m} نحو
\qty{20}{kyr}، وهذا متّسق.
\end{solution}

\begin{solution}{exo:b3:sequence-stratigraphy:5}
$\mathrm dx_s/\mathrm dt = (4000 - 0.15\times20\,000)/50 = \qty{20}{m/kyr}$،
نحو البحر. يبقى خط الشاطئ ساكنًا عندما $AW = Q$: $A = 4000/20\,000 =
\qty{0.2}{m/kyr}$.
\end{solution}

\begin{solution}{exo:b3:sequence-stratigraphy:6}
$\mathrm dx_s/\mathrm dt = 4000/50 + 0.2/0.01 = 80 + 20 = \qty{100}{m/kyr}$؛
يفسّر الانخفاض وحده \qty{20}{m/kyr}، أي خُمسها.
\end{solution}

\begin{solution}{exo:b3:sequence-stratigraphy:7}
يتراجع الخط عندما $Q < AW = 0.4\times25\,000 = \qty{10000}{m^3/kyr}$ لكلّ
متر. ومع $Q = 6000$: $(6000 - 10\,000)/10 = \qty{-400}{m/kyr}$، أي تراجعًا
بمقدار \qty{0.4}{km} كلّ ألف سنة: الوجه الشاطئي المنخفض يُنتج
تجاوزًا بحريًّا سريعًا.
\end{solution}

\begin{solution}{exo:b3:sequence-stratigraphy:8}
$1/94.9 - 1/123.9 = \num{0.00247}\text{ دورة}/\unit{kyr}$: يلزم سجلّ أطول من
$1/0.00247 = \qty{406}{kyr}$ يفصل بينهما، و\qty{600}{kyr} (الدقّة
0.00167)، وهذا يكفي. $1/94.9 - 1/98.9 = 0.00042$: يحتاج الأمر نحو
\qty{2400}{kyr}، أي أربعة أضعاف طول السجلّ.
\end{solution}

\begin{solution}{exo:b3:sequence-stratigraphy:9}
$16.4/(20\times40.98) = \qty{0.0200}{m/kyr}$، \qty{2.0}{cm/kyr}؛ طبقات الرماد
تفصل بينها $3.1/0.020 = \qty{155}{kyr}$.
\end{solution}

\begin{solution}{exo:b3:sequence-stratigraphy:10}
(أ) $\tan\alpha = -0.4/200 = -0.002$: هبوط بزاوية \ang{0.11}. (ب)
$0.2/50 = 0.004$: صعود بزاوية \ang{0.23}. يمتدّ المقطع الزلزالي عشرات
الكيلومترات بتضخيم رأسي من عشرات إلى مئات، لذلك تظهر الزاوية؛
أمّا جرف طوله بضع مئات من الأمتار فيُظهر ارتفاعًا لا يتجاوز مترًا، يضيع في
تغيّر الطبقات.
\end{solution}

\begin{solution}{exo:b3:sequence-stratigraphy:11}
$1\,296\,000/54.725282 = \qty{23.68}{kyr}$ ($g_5$)؛ $/57.924383 =
\qty{22.37}{kyr}$ ($g_2$)؛ $/68.377912 = \qty{18.95}{kyr}$ ($g_4$)؛
$/67.834313 = \qty{19.11}{kyr}$ ($g_3$). $g_2 - g_5 = 3.199101''$ لكلّ سنة،
\qty{405.1}{kyr}؛ $g_4 - g_5 = 13.65263''$، \qty{94.9}{kyr}.
\end{solution}

\begin{solution}{exo:b3:sequence-stratigraphy:12}
يُنتجان كلاهما $0.3\times100 = \qty{30}{m}$ من حيّز الترسيب بالمعدّل نفسه؛ ومع
التغذية نفسها تتطابق الهندسة، لأنّ الرواسب تستجيب فقط لمستوى
البحر النسبي. التغيّر الأيوستاتي واحد في كلّ مكان، فالتتابع
نفسه، بالعمر نفسه، على هوامش أخرى يشير إلى التاريخ الأول. يُبنى تاريخ الهبوط
من تاريخ الحرارة والحمولة في الحوض
(إزالة الطبقات)، أو سجلّ أيوستاتي مستقلّ كنظائر الأكسجين،
ويفصل بينهما أيضًا.
\end{solution}

\begin{solution}{pb:b3:sequence-stratigraphy:1}
\textbf{1.} $\dot E = 30\times2\pi/400\times\sin(2\pi t/400)$: على الأكثر
\qty{0.471}{m/kyr}، ويرتفع عند \qty{100}{kyr} و\qty{500}{kyr}، وينخفض عند
\qty{300}{kyr}.

\textbf{2.} \qty{0.020}{m/kyr} عند المفصل، و0.037 عند \qty{20}{km}، و0.062 عند
\qty{50}{km}، و\qty{0.120}{m/kyr} عند \qty{120}{km}.

\textbf{3.} $A < 0$ عندما $\sin(2\pi t/400) < -0.0346/0.471 = -0.0734$: $t =
200 + (400/2\pi)\arcsin0.0734 = \qty{204.7}{kyr}$.

\textbf{4.} مع $\sigma = 0.0617$: $t = 400 - 63.66\arcsin0.131 =
\qty{391.6}{kyr}$.

\textbf{5.} نحو \qty{187}{kyr} مقابل 200: يجب أن يتغلّب الانخفاض الأيوستاتي أوّلًا
على الهبوط، ويكفّ عن تغلّبه عليه قبل بلوغ المستوى المنخفض.

\textbf{6.} $0.045\times400 = \qty{18}{m}$، مقابل مدى أيوستاتي قدره
\qty{60}{m}: أي الثلث.

\textbf{7.} نحو اليابسة من 50 حتى نحو \qty{120}{kyr}، ونحو البحر حتى نحو
\qty{440}{kyr}، ونحو اليابسة حتى نحو \qty{560}{kyr}، ونحو البحر مجدّدًا حتى
\qty{600}{kyr} (نقاط الانعطاف في النموذج هي 118.5 و439.5
و\qty{559}{kyr}؛ جدول كلّ \qty{50}{kyr} لا يحصرها إلّا بين حدّين).

\textbf{8.} $(6000 - 0.30\times15\,000)/40 = \qty{37.5}{m/kyr}$، نحو البحر.

\textbf{9.} $6000/15\,000 = \qty{0.40}{m/kyr}$.

\textbf{10.} $6000/35.7 + 0.424/0.012 = 168.1 + 35.3 = \qty{203}{m/kyr}$,
ضمن \qty{0.5}{\percent} من النموذج.

\textbf{11.} التغذية: \qty{83}{\percent} من السرعة.

\textbf{12.} $(-42.7 - 16.3)/(50.7 - 17.4) = \qty{-1.77}{m/km}$, \ang{0.10}.
منذ ذلك الحين هبطت النقطة الأولى بمقدار $0.0345\times395.5 =
\qty{13.6}{m}$، والثانية $0.0623\times208.5 = \qty{13.0}{m}$: المقداران متقاربان،
لذلك يختلف الهبوط عند المناسيب الحالية، \qty{59.0}{m}، عن الانخفاض الأيوستاتي
بمقدار \qty{0.6}{m} فقط. ويكون الانخفاض النسبي عند خط الشاطئ في ذلك الوقت،
أي الأيوستاسيا مطروحًا منها الهبوط خلال الانخفاض، أصغر.

\textbf{13.} $150$ و200: المستوى المرتفع؛ 250 و300 و350: طور الانخفاض؛ 400:
المستوى المنخفض؛ 450 و500 و550: تجاوزي (يتراجع خط الشاطئ مباشرةً بعد
440)؛ 600: المستوى المرتفع.

\textbf{14.} $204.7$ و\qty{391.6}{kyr} مقابل 204.5 و391.5: متطابقان،
ضمن خطوة النموذج البالغة \qty{0.5}{kyr}.

\textbf{15.} على الرصيف، السطح المكشوف أثناء الانخفاض: قمّة
رسوبيات المستوى المرتفع من المفصل حتى \qty{17}{km}، وتتابع الامتداد فوق
الواجهات الأمامية المتدرّجة الهبوط لطور الانخفاض؛ وفي الحوض، التوافق المقابل
عند قاعدة مسلك طور الانخفاض أو قمّته، وفق الاصطلاح. والرصيف
الحقيقي سيُقطع بالأنهار إلى أودية محفورة، وتتعرّض للتجوية فتتحوّل إلى تربة،
ويتعرّض جزء منه للتعرية.

\textbf{16.} $(18.5 - 52.8)/(559 - 439.5) = \qty{-0.29}{km/kyr}$؛ ويتكوّن
\omterm{def:b3:sequence-stratigraphy:surfaces}{سطح الغمر الأقصى} في نهايته.

\textbf{17.} يجتاح خط الشاطئ بسرعة سهلًا ساحليًا شبه مستوٍ:
تُحجَز التغذية قربه، وتنتشر فوق مساحة واسعة، ويتلقّى كلّ موضع
قدرًا ضئيلًا منها. ونحو الحوض انطلاقًا من قدم المنحدر، بعد نحو \qty{55}{km}، يكون
قاع البحر مفتقرًا إلى الرواسب، فيتشكّل هناك \omterm{def:b3:sequence-stratigraphy:condensed}{المقطع المكثَّف}.

\textbf{18.} الاستبانة $1/30 = \num{0.033}\text{ دورة}/\unit{m}$؛ نايكويست $1/(2\times
0.01) = \num{50}\text{ دورة}/\unit{m}$.

\textbf{19.} $7.5/2.0 = 3.75$، 1 و$0.81/2.0 = 0.405$؛ النسب المدارية
هي $405/100 = 4.05$ و$41/100 = 0.41$.

\textbf{20.} \qty{2.0}{m} لكلّ \qty{100}{kyr}: \qty{2.0}{cm/kyr}؛ $0.81/41 =
\qty{1.98}{cm/kyr}$؛ $7.5/405 = \qty{1.85}{cm/kyr}$. أمّا أطوال الموجات التي
يستطيع \omterm{def:b3:sequence-stratigraphy:periodogram}{المخطط الدوري} إظهارها قرب \qty{8}{m} هي 10 و7.5 و\qty{6}{m}: والقيمة
الحقيقية \qty{8.1}{m} تقع بين اثنتين منها.

\textbf{21.} $30/0.020 = \qty{1500}{kyr}$.

\textbf{22.} $23.7\times0.020 = \qty{0.474}{m}$. عند أخذ عيّنة كلّ
\qty{25}{cm}، يكون تردّد نايكويست 2 دورة لكلّ متر، أي طول موجة قدره
\qty{0.5}{m}: إنّ التقدّم الاعتدالي، عند 2.11 دورة لكلّ متر، يقع خارجه وينعكس
إلى $4 - 2.11 = 1.89$، فيعطي طول موجة كاذبًا قدره \qty{0.53}{m}.

\textbf{23.} قمّة واحدة تُطابق دورةً ما عند معدّلٍ ما، أيًّا كان
مصدرها؛ أمّا ثلاث قمم تُطابق نسبُها النسب المدارية فهي مستقلّة عن
المعدّل، الذي يجب أن يخرج حينئذٍ متطابقًا من كلّ منها.

\textbf{24.} ينتقل خط الشاطئ من \qty{17.37}{km} عند بداية
الانخفاض حتى \qty{52.84}{km} عند نهاية انحسار المستوى المنخفض: أقصى
إزاحة نحو الحوض بمقدار \textbf{\qty[text-series-to-math=true]{35.5}{km}}.
\end{solution}
